\BourbakiExerciseGroup

\begin{enumerate}
\def\labelenumi{\arabic{enumi})}
\item
  Let \(\mathbf{K}\) be a field of characteristic \(p\) , \(n\) an integer not a multiple of \(p\) , \(\mathbf{R}_n(\mathbf{K})\) the field of \(n\) -th roots of unity and \(\mathbf{E}\) a cyclic extension of \(\mathbf{R}_n(\mathbf{K})\) of degree \(n\) over \(\mathbf{K}\) ; \(\mathbf{E}\) is thus the root field of a polynomial \(X^{n}-\alpha\in\mathbb{R}_{n}(\mathbf{K})[\mathbf{X}]\) (V, p.~89, Th. 4). For E to be a Galois extension of K it is necessary and sufficient that for every automorphism \(\tau\in\operatorname{Gal}\left(\mathbb{R}_{n}(\mathbf{K})/\mathbf{K}\right)\) there exists an integer r\textgreater0 and an element \(\gamma\in\mathbb{R}_{n}(\mathbf{K})\) such that \(\tau(\alpha)=\gamma^{n}\alpha^{r}\) .
\item
  \begin{enumerate}
  \def\labelenumii{\alph{enumii})}
  \tightlist
  \item
    Let E be a splitting extension of the polynomial \(X^{3}-2\in Q[X]\) . The field E is a cyclic extension of degree 3 of \(\mathrm{R}_{3}(\mathbf{Q})\) but contains no cyclic extension of degree 3 of Q. b) The field \(\mathrm{E}=\mathrm{R}_{9}(\mathbf{Q})\) is a cyclic extension of degree 3 of \(\mathrm{R}_{3}(\mathbf{Q})\) and contains a cyclic extension of degree 3 of Q.
  \end{enumerate}
\end{enumerate}

\begin{enumerate}
\def\labelenumi{\alph{enumi})}
\setcounter{enumi}{2}
\tightlist
\item
  With the hypothesis and notation as in V, p.~154, Ex. 7, show that for \(\mathrm{K}(\alpha,\beta)\) to be Galois over K it is necessary and sufficient for d to be a square in \(\mathrm{K}(\alpha)\) ; for \(\mathrm{K}(\beta)\) to be cyclic over K it is necessary and sufficient for d to be a square in K.
\end{enumerate}

* 3) Show that for every integer \(n > 1\) there exists a cyclic extension of degree \(n\) of the field \(\mathbf{Q}\) . (Reduce to the case where \(n\) is a prime number \(q\) and use the structure of the multiplicative group \((\mathbf{Z} / q^m\mathbf{Z})^*\) (VII, p.~13). *

\begin{enumerate}
\def\labelenumi{\arabic{enumi})}
\setcounter{enumi}{3}
\tightlist
\item
  Let \(\mathbf{K}\) be a field of characteristic \(p, q\) a prime number \(\neq p\) such that \(\mathbf{K}\) contains the \(q\) -th roots of unity and let \(\zeta\) be a primitive \(q\) -th root of unity.
\end{enumerate}

\begin{enumerate}
\def\labelenumi{\alph{enumi})}
\item
  Let E be a cyclic extension of K of degree \(q^e\) and \(\sigma\) a K-automorphism of E generating Gal(E/K). Let F be the field of degree \(q^{e-1}\) over K such that \(K \subset F \subset E\) ; there exists \(\theta \in E\) , a root of an irreducible polynomial \(X^q - \alpha\) in F{[}X{]} such that \(E = F(\theta)\) and \(\theta^{\sigma^m} = \zeta\theta\) (where \(m = q^{e-1}\) ). Show that also \(E = K(\theta)\) and \(\theta^\sigma = \beta\theta\) where \(\beta \in F\) is such that \(\alpha^{\sigma-1} = \beta^q\) and \(N_{F/K}(\beta) = \zeta\) (observe that \(\theta^\sigma = \beta\theta^k\) with \(0 < k < q\) and \(\beta \in F\) , by Ex. 1; by calculating \(\theta^{\sigma^{mq}}\) , show that \(k^{mq} - 1 \equiv 0 \pmod{q}\) ) and deduce that \(k = 1\) ).
\item
  Conversely, let F be a cyclic extension of degree \(q^{e-1}\) of K with \(e \geqslant 2\) and let \(\sigma\) be a K-automorphism of F generating Gal(F/K). If there exists an element \(\beta \in F\) such that \(N_{F/K}(\beta) = \zeta\) and if \(\alpha \in F\) is such that \(\alpha^{\sigma-1} = \beta^q\) , show that for every \(c \in K^*\) the polynomial \(X^q - c\alpha\) is irreducible in F{[}X{]}. If \(\theta\) is one of the roots of this polynomial in an algebraic closure \(\Omega\) of F, show that there exists a K-homomorphism \(\bar{\sigma}\) of \(E = F(\theta)\) into \(\Omega\) , extending \(\sigma\) and such that \(\theta^{\sigma} = \beta\theta\) ; deduce that E is a cyclic extension of K of degree \(q^e\) , that \(\bar{\sigma}\) generates Gal(E/K) and that \(E = K(\theta)\) . Finally show that every cyclic extension of degree \(q^e\) of K, containing F, is the root field of a polynomial \(X^q - c\alpha\) in F{[}X{]}, for an appropriate \(c \in K^*\) (use Th. 3 of V, p.~85).
\item
  Take K to be the field Q of rational numbers. The polynomial \(X^{2} + 1\) is irreducible in Q{[}X{]}; if i is one of its roots, then \(F = \mathbf{Q}(i)\) is a cyclic extension of Q of degree 2 but there exists no cyclic extension of degree 4 of Q containing F.
\end{enumerate}

\begin{enumerate}
\def\labelenumi{\arabic{enumi})}
\setcounter{enumi}{4}
\item
  Let K be a field of characteristic p and n an integer not a multiple of p.~For a polynomial of K{[}X{]} of the form \(X^{n} - a\) to be irreducible it is necessary and sufficient that for every prime factor q of n, a should not equal the q-th power of an element of K and further when \(n \equiv 0 \pmod{4}\) , a should not be of the form \(-4c^{4}\) with \(c \in K\) . (To prove the sufficiency of the condition, reduce to the case where \(n = q^{e}\) with q prime and use Ex. 4 of V, p.~153. Argue by induction an e, determining, with the help of Ex. 4 of V, p.~153, the form of the constant term of each irreducible factor of \(X^{q^{e}} - a\) .)
\item
  Let \(\mathbf{K}\) be a field of characteristic \(p, n\) an integer not a multiple of \(p\) and \(\mathbf{N}\) a splitting field of the polynomial \(\mathbf{X}^n - a\) of \(\mathbf{K}[\mathbf{X}]\) .
\end{enumerate}

\begin{enumerate}
\def\labelenumi{\alph{enumi})}
\tightlist
\item
  Show that the group \(\operatorname{Gal}(\mathbf{N} / \mathbf{K})\) is isomorphic to a subgroup of the group \(\Gamma\) of matrices \(\left( \begin{array}{cc}x & y\\ 0 & 1 \end{array} \right)\) , where \(x\in (\mathbf{Z} / n\mathbf{Z})^*\) and \(y\in \mathbf{Z} / n\mathbf{Z}\) . If \(n\) is a prime number, the group
\end{enumerate}

Gal(N/K) is commutative only if \(\mathrm{R}_{n}(\mathrm{K})=\mathrm{K}\) or if a is the n-th power of an element of K. b) If K=Q and if n is a prime number, show that if a is not the n-th power of a rational number, then Gal(N/K) is isomorphic to the whole of \(\Gamma\) .

\begin{enumerate}
\def\labelenumi{\alph{enumi})}
\setcounter{enumi}{2}
\tightlist
\item
  If \(\mathbf{K} = \mathbf{Q}\) , determine \(\operatorname{Gal}(\mathbf{N} / \mathbf{K})\) when \(a\) is a square-free integer.
\end{enumerate}

\begin{enumerate}
\def\labelenumi{\arabic{enumi})}
\setcounter{enumi}{6}
\tightlist
\item
  An extension E of a field K is called a Kummer extension if there exists an integer n \textgreater{} 0 such that K contains a primitive n-th root of unity and E is an abelian extension with Galois group annihilated by n.
\end{enumerate}

\begin{enumerate}
\def\labelenumi{\alph{enumi})}
\item
  Let K be a field and N a Galois extension of degree m of K with Galois group \(\Gamma\) . For N to be a Kummer extension of K it is necessary and sufficient that there should exist a basis \((\theta_{\sigma})_{\sigma\in\Gamma}\) of N over K such that the elements \(\gamma_{\sigma\tau}=\theta_{\sigma}^{1-\tau}\) belong to K for each pair \((\sigma,\tau)\) . (For the necessity of the condition reduce to the case where N is cyclic of prime power degree. For the sufficiency observe that if \(\tau\in\Gamma\) is of order n, then there exists \(\sigma\in\Gamma\) such that \(\gamma_{\sigma\tau}\) is a primitive n-th root of unity.)
\item
  If \((\theta_{\sigma})_{\sigma \in \Gamma}\) is a basis satisfying condition \(a)\) , show that for an element \(x = \sum_{\sigma} a_{\sigma} \theta_{\sigma}\) of \(N\) to
\end{enumerate}

be such that the conjugates of \(x\) form a normal basis of \(\mathbf{N}\) over \(\mathbf{K}\) , it is necessary and sufficient that \(a_{\sigma} \neq 0\) for all \(\sigma \in \Gamma\) .

* 8) Let N be a Kummer extension of a field K.

\begin{enumerate}
\def\labelenumi{\alph{enumi})}
\item
  Let E be a subextension of N such that \([E:K]=q\) is a prime number. Show that there exists a decomposition \(N=N_{1}\otimes N_{2}\otimes\cdots\otimes N_{r}\) of N into cyclic extensions of K such that \(E\subset N_{j}\) for an index j.
\item
  Deduce from a) that if \(x \in N\) is such that the conjugates of x over K form a normal basis of N over K, then for each field E such that \(K \subset E \subset N\) , the conjugates of x over E form a normal basis of N over E (reduce to the case where \([E : K]\) is prime and use Ex. 7, b)). *
\end{enumerate}

\begin{enumerate}
\def\labelenumi{\arabic{enumi})}
\setcounter{enumi}{8}
\tightlist
\item
  Let \(\mathbf{K}\) be a field of characteristic \(p > 0\) .
\end{enumerate}

\begin{enumerate}
\def\labelenumi{\alph{enumi})}
\item
  Let E be a cyclic extension of degree \(p^e\) of K and let \(\sigma\) be a K-automorphism generating the group \(\operatorname{Gal}(\mathrm{E} / \mathrm{K})\) . Let F be the field of degree \(p^{e - 1}\) over K such that \(\mathbf{K} \subset \mathbf{F} \subset \mathbf{E}\) ; if \(m = p^{e - 1}\) , show that there exists \(\theta \in \mathbf{E}\) , root of an irreducible polynomial \(\mathbf{X}^p - \mathbf{X} - \alpha\) of \(\mathrm{F}[\mathbf{X}]\) such that \(\mathrm{E} = \mathrm{F}(\theta)\) and \(\sigma^m (\theta) = \theta + 1\) . Show that we also have \(\mathrm{E} = \mathrm{K}(\theta)\) and \(\sigma (\theta) = \theta + \beta\) where \(\beta \in \mathrm{F}\) is such that \(\sigma (\alpha) - \alpha = \beta^p - \beta\) and \(\mathrm{Tr}_{\mathrm{F} / \mathrm{K}}(\beta) = 1\) .
\item
  Conversely, let F be a cyclic extension of degree \(p^{e-1}\) of K (with e \textgreater{} 1), \(\sigma\) a K-automorphism of F generating Gal(F/K). Show that there exist two elements \(\alpha\) , \(\beta\) of F such that \(\operatorname{Tr}_{\mathrm{F}/\mathrm{K}}(\beta) = 1\) and \(\sigma(\alpha) - \alpha = \beta^{p} - \beta\) (use Th. 3 of V, p.~85 and the Cor. of Prop. 1 of V, p.~49). Deduce that for all \(c \in K\) the polynomial \(X^{p} - X - \alpha - c\) is irreducible in F{[}X{]}; if \(\theta\) is a root of this polynomial in an algebraic closure \(\Omega\) of F, show that there exists a K-homomorphism \(\bar{\sigma}\) of E = F( \(\theta\) ) in \(\Omega\) extending \(\sigma\) and such that \(\bar{\sigma}(\theta) = \theta + \beta\) ; conclude that E is a cyclic extension of degree \(p^{e}\) of K, that \(\bar{\sigma}\) generates Gal(E/K) and that E = K( \(\theta\) ). Finally show that every cyclic extension of degree \(p^{e}\) of K containing F is the splitting field of a polynomial \(X^{p} - X - \alpha - c\) of F{[}X{]} for a suitable \(c \in K\) .
\end{enumerate}

\begin{enumerate}
\def\labelenumi{\arabic{enumi})}
\setcounter{enumi}{9}
\tightlist
\item
  \begin{enumerate}
  \def\labelenumii{\alph{enumii})}
  \tightlist
  \item
    Let \(\mathbf{K}\) be a field whose algebraic closure is an extension of prime degree \(q\) of \(\mathbf{K}\) . Such that \(\mathbf{K}\) is perfect and that \(q\) must be distinct from the characteristic of \(\mathbf{K}\) (use Ex. 9).
  \end{enumerate}
\end{enumerate}

\begin{enumerate}
\def\labelenumi{\alph{enumi})}
\setcounter{enumi}{1}
\tightlist
\item
  Show that K contains the q-th roots of unity and that \(\Omega\) is the splitting field of an irreducible polynomial \(X^{q}-a\) of K{[}X{]}; deduce that q=2 (in the contrary case deduce from Ex. 5 that the polynomial \(X^{q^{2}}-a\) would be irreducible). Show further that -a is a square in K (loc. cit.), that -1 is not a square in K, and deduce that \(\Omega=\mathrm{K}(i)\) , where \(i^{2}=-1\) .
\end{enumerate}

\begin{enumerate}
\def\labelenumi{\arabic{enumi})}
\setcounter{enumi}{10}
\item
  Let \(\mathbf{K}\) be a field whose algebraic closure \(\Omega\) is an extension of finite degree \(> 1\) of \(\mathbf{K}\) . Show that \(\Omega = \mathbf{K}(i)\) where \(i^2 = -1\) . (If we had \(\Omega \neq \mathbf{K}(i)\) , show with the help of Galois theory and Sylow's theorem that there would then exist a field \(E\) such that \(\mathbf{K}(i) \subset E \subset \Omega\) and \(\Omega\) is an extension of prime degree of \(E\) ; now apply Ex. 10.)
\item
  Let K be a field of characteristic q, \(\Omega\) an algebraically closed extension of K, x an element of \(\Omega\) such that \(x \notin K\) and M a maximal element of the set of subextensions of \(\Omega\) not containing x (V, p.~156, Ex. 9); suppose further that M is perfect so that \(\mathbf{M}(x)\) is a Galois extension of prime degree p of M (loc. cit.).
\end{enumerate}

\begin{enumerate}
\def\labelenumi{\alph{enumi})}
\item
  Show that either \(p = 2\) and \(\Omega = M(i)\) or for each integer \(r \geqslant 1\) , there exists a single extension of degree \(p^r\) of \(M\) ; this extension \(M_r\) is cyclic over \(M\) , \(\Omega\) is the union of the \(M_r\) and the \(M_r\) are the only extensions of finite degree of \(M\) . (Note that if a \(p\) -group \(\Gamma\) of order \(p^r\) has only one subgroup of order \(p^{r-1}\) , it is necessarily cyclic, arguing by induction on \(r\) and using I, p.~133, Ex. 10, and I, p.~77, Cor. of Prop. 11).
\item
  Suppose that \(\Omega \neq \mathbf{M}(i)\) and that \(p \neq q\) . Show that \(\mathbf{M}\) contains the \(p\) -th roots of unity; we thus have \(\mathbf{M}(x) = \mathbf{M}_1 = \mathbf{M}(\alpha)\) with \(\alpha^p \in \mathbf{M}\) . Show that if \(\alpha\) is the \(p\) -th power of an element of \(\mathbf{M}_1\) , then \(\mathbf{M}_1 = \mathbf{M}(\zeta)\) , where \(\zeta\) is \(p^2\) -th root of unity (if \(\alpha = \beta^p\) for \(\beta \in \mathbf{M}_1\) , consider the minimal polynomial of \(\beta\) in \(\mathbf{M}[\mathbf{X}]\) ).
\item
  Under the hypotheses of \(b\) ) show that if \(\alpha\) is not a \(p\) -th power of an element of \(\mathbf{M}_1\) then \(\mathbf{M}_r\) is the splitting field of the polynomial \(\mathbf{X}^{p^{r-1}} - \alpha\) of \(\mathbf{M}_1[\mathbf{X}]\) (use \(b\) )).
\end{enumerate}

\(d)\) Under the hypotheses of \(b\) ), show that if \(\alpha\) is the \(p\) -th power of an element of \(\mathbf{M}_1\) , there exists \(\gamma \in \mathbf{M}_1\) such that \(\gamma \notin \mathbf{M}\) and that \(\mathbf{M}_2\) is the splitting field of \(\mathbf{X}^p - \gamma\) ; show that \(\gamma\) is not the \(p\) -th power of an element of \(\mathbf{M}_2\) and deduce that \(\mathbf{M}_r\) is the splitting field of the polynomial \(\mathbf{X}^{p^{r-1}} - \gamma\) of \(\mathbf{M}_1[\mathbf{X}]\) .

\begin{enumerate}
\def\labelenumi{\arabic{enumi})}
\setcounter{enumi}{12}
\tightlist
\item
  A field K is called a Moriya field if every separable algebraic extension of finite degree over K is cyclic. * Every finite field is a Moriya field (V, p.~95, Prop. 4). * The field M of Ex. 12 is a Moriya field.
\end{enumerate}

\begin{enumerate}
\def\labelenumi{\alph{enumi})}
\item
  For K to be a Moriya field it is necessary and sufficient that for each integer n \textgreater{} 0 there should exist at most one separable extension of degree n over K. (Use the following fact: if G is a finite group and for each integer \(m \geqslant 1\) dividing the order of G there exists at most one subgroup of G of order m, then G is cyclic; this follows from the same property for p-groups (Ex. 12) and from I, p.~80, Th. 4).
\item
  If K is a Moriya field, show that every algebraic extension E of K is a Moriya field and that if F is an extension of E of degree m over E, then there exists an extension of K of degree m over K.
\end{enumerate}

\begin{enumerate}
\def\labelenumi{\arabic{enumi})}
\setcounter{enumi}{13}
\tightlist
\item
  \begin{enumerate}
  \def\labelenumii{\alph{enumii})}
  \tightlist
  \item
    Let K be a field such that there exist algebraic extensions of K of arbitrarily large finite degree, and that all these extensions have as degree a power of the same prime number p.~Show that under these conditions, for each power \(p^{h}\) ( \(h \geqslant 0\) ) there exists an extension of K of degree \(p^{h}\) (use I, p.~77, Th. 1).
  \end{enumerate}
\end{enumerate}

\begin{enumerate}
\def\labelenumi{\alph{enumi})}
\setcounter{enumi}{1}
\tightlist
\item
  Let p be a prime number and K a field such that there exists an algebraic extension of K of degree divisible by p if \(p \neq 2\) and divisible by 4 if p = 2. Show that there exists an algebraic extension E of K such that the set of degrees of algebraic extensions of E is the set of powers \(p^{h}\) of p.~(If K is perfect or of characteristic \(\neq p\) , let \(K_{1}\) be the perfect closure of K; consider a maximal algebraic extension of \(K_{1}\) , among all those whose elements are of degree not divisible by p over \(K_{1}\) ; use a) and Ex. 11.)
\end{enumerate}

\begin{enumerate}
\def\labelenumi{\arabic{enumi})}
\setcounter{enumi}{14}
\item
  Let K be a Moriya field (Ex. 13) and P the set of prime numbers dividing the degree of an algebraic extension of finite degree of K. Show that the set of degrees of algebraic extensions of finite degree of K is either equal to the set \(N_{P}\) of integers all of whose prime factors are in P or to the subset of \(N_{P}\) consisting of integers not divisible by 4.
\item
  For every prime number \(p\) there exists a Moriya field of characteristic \(p\) such that the set of degrees of algebraic extensions of finite degree of this field is the set \(\mathbf{N}^*\) of integers \(>0^1\) . Show that for every set \(P\) of prime numbers, there exists a Moriya field \(K\) of any characteristic, such that the set of degrees of algebraic extensions of finite degree of \(K\) is the set \(N_P\) (argue as in Ex. 14, b)).
\end{enumerate}

* 17) Let \(\mathbf{P}_n\) be the set of monic polynomials \(\mathbf{X}^n + a_1\mathbf{X}^{n-1} + \cdots + a_n\) of \(\mathbf{Z}[\mathbf{X}]\) whose roots in \(\mathbf{C}\) all have absolute value \(\leqslant 1\) .

\begin{enumerate}
\def\labelenumi{\alph{enumi})}
\item
  Show that the set \(P_{n}\) is finite. Deduce that for every polynomial \(F \in P_{n}\) , if \(F_{h}\) denotes the polynomial whose roots are the h-th powers of the roots of F, then there exist two distinct integers h, k such that \(F_{h} = F_{k}\) .
\item
  Conclude from a) that all the roots of a polynomial \(F \in P_{n}\) are either zero or roots of unity (« Kroneckers's theorem »). *
\end{enumerate}

\begin{enumerate}
\def\labelenumi{\arabic{enumi})}
\setcounter{enumi}{17}
\tightlist
\item
  Let \(p \in \mathbb{N}\) be a prime number, \(C_p\) the cyclotomic extension of level \(p\) of \(\mathbf{Q}\) , \(\zeta \in C_p\) a primitive \(p\) -th root of unity and \(n\) an integer. Write
\end{enumerate}

\[
\mathrm{X} _ {n} (\zeta) = \sum_ {\alpha \in \mathbf {Z} / p \mathbf {Z}} \zeta^ {\alpha^ {n}}.
\]

\begin{enumerate}
\def\labelenumi{\alph{enumi})}
\item
  Calculate \(d = [\mathbf{Q}(\mathbf{X}_n(\zeta)) : \mathbf{Q}]\) . Determine the Galois group of the extension \(\mathbf{Q}(\mathbf{X}_n(\zeta))\) of \(\mathbf{Q}\) and the conjugates of \(\mathbf{X}_n(\zeta)\) .
\item
  Let \(T^{d} + a_{1}T^{d-1} + \cdots + a_{d}\) be the minimal polynomial of \(\mathbf{X}_{n}(\zeta)\) over Q. Show that \(a_{1} = 0\) .
\item
  Suppose that \(n > 1\) and \(p \equiv 1 \pmod{n}\) . Show that \(a_2 = -\frac{n(n - 1)p}{2}\) if \(\frac{p - 1}{n}\) is even and \(a_2 = \frac{np}{2}\) if \(\frac{p - 1}{n}\) is odd.
\item
  Suppose that \(p \equiv 1 \pmod{n}\) . Let \(\beta \in \mathbf{F}_p^*\) be an element whose residue class in \(\mathbf{F}_p^* / \mathbf{F}_p^{*n}\) is a generator. Let \(\mu\) be the number of solutions in \(\mathbf{F}_p^n\) of the equation \(\sum_{i=1}^{n} \beta^{i-1} \mathbf{X}_i^n = 0\) . Show that \(a_n = \frac{(-1)^n p(\mu - p^{n-1})}{p-1}\) .
\item
  Calculate \(a_3\) when \(n = 3\) and \(p = 7, 13, 19\) .
\end{enumerate}

*1 The largest totally ramified extension of the field \(\mathbf{Q}_p\) of \(p\) -adic numbers in its algebraic closure is a field of characteristic 0 with the same property. *

* 19) a) Assume the prime number theorem; if \(\pi(x)\) is the number of primes \(p \leqslant x\) , then \(\pi(x) \sim x / \log x\) as \(x\) tends to \(+\infty\) . Deduce that for every odd integer \(m\) , there exists an increasing sequence of prime numbers \(p_1 < p_2 < \cdots < p_m\) such that \(p_m < p_1 + p_2\) (consider the prime numbers between \(x\) and \(3x/2\) ).

\begin{enumerate}
\def\labelenumi{\alph{enumi})}
\setcounter{enumi}{1}
\tightlist
\item
  With the finite sequence of prime numbers \(p_{j}\) ( \(1 \leq j \leq m\) ) determined as in a), put \(n = p_{1}p_{2} \ldots p_{m}\) . Show that in the cyclotomic polynomial \(\Phi_{n}(X)\) the terms of degree \(< p_{m} + 1\) are the same as those of the polynomial
\end{enumerate}

\[
(1 - \mathbf {X}) ^ {- 1} \prod_ {j = 1} ^ {m} (1 - \mathbf {X} ^ {p _ {j}}).
\]

Deduce that the coefficient of \(\mathbf{X}^{p_m}\) in \(\Phi_n(\mathbf{X})\) is \(1 - m\) (« I. Schur's theorem »). *

* 20) Let \(p\) be a prime number and \(G\) a commutative \(p\) -torsion group (VII, p.~10).\\
a) Let \(G_i \subset G\) be the set of elements \(x \in G\) such that for each \(l\) the equation \(p^{l} y = x\) has a solution. Show that \(G_i\) is a direct factor of \(G\) .

\begin{enumerate}
\def\labelenumi{\alph{enumi})}
\setcounter{enumi}{1}
\item
  Suppose that the equation px = 0 has only a finite number of solutions in G. Show that for each integer l, the kernel \(_{l}G\) of the multiplication by \(p^{l}\) in G is a finite group and that the rank \(r(l)\) of the vector \(F_{p}\) -space \(_{l}G/_{(l-1)}G\) is a decreasing function of l, hence constant for \(l > l_{0}\) , \(l_{0}\) large enough. (Use the theorem of elementary divisors (VII, p.~24). Show that \(G/G_{i}\) is a finite group annihilated by \(p^{l_{0}}\) .)
\item
  Suppose further that \(\mathbf{G} = \mathbf{G}_i\) . Show that the mapping \(l \mapsto r(l)\) is constant. Writing \(\mathbf{H} = (\mathbf{Z}[1/p]/\mathbf{Z})^{(1)}\) , show that for each \(l\) there exists an isomorphism \(_l\mathbf{G} \to {}_l\mathbf{H}\) and that every isomorphism \(_l\mathbf{G} \to {}_l\mathbf{H}\) extends to an isomorphism \((l+1)\mathbf{G} \to {}_{(l+1)}\mathbf{H}\) . Deduce that \(\mathbf{G}\) is isomorphic to \(\mathbf{H}\) .
\item
  Under the hypotheses of b) show that there exists an integer n and an application \(i \to n_i\) of \(N - \{0\}\) into N, with finite support, such that G is isomorphic to \((\mathbf{Z}[1/p]/\mathbf{Z})^n \oplus \bigoplus_{i} (\mathbf{Z}/p^i \mathbf{Z})^{n_i}\) . Show that the integer n and the mapping \(i \mapsto n_i\) are uniquely
\end{enumerate}

determined by the mapping \(l \mapsto \text{Card}(l\mathbf{G})\) . *

* 21) Let G and H be two commutative torsion groups such that for every integer \(n\) , the equation \(nx = 0\) has the same finite number of solutions in G and in H. Show that G is isomorphic to H. (Decompose G and H into direct sums of \(p\) -torsion groups and apply Ex. 20.) Let \(G = \bigoplus_{i \in \mathbf{N}} \mathbf{Z} / p^i\mathbf{Z}\) and \(H = \bigoplus_{i \in \mathbf{N}} \mathbf{Z} / p^{2i}\mathbf{Z}\) . Show that G and H are not isomorphic

and that for each integer \(n\) the cardinal of the set of solutions of the equation \(nx = 0\) is the same in \(G\) and in \(H\) . *

\begin{enumerate}
\def\labelenumi{\arabic{enumi})}
\setcounter{enumi}{21}
\tightlist
\item
  \begin{enumerate}
  \def\labelenumii{\alph{enumii})}
  \tightlist
  \item
    Let k be a field, f a separable monic polynomial in \(k[X]\) , D a splitting extension of \(f, x_{1}, \ldots, x_{n}\) the roots of f in D and G the Galois group of D over k. Assume that G is cyclic. Show that for every element of G to induce an even permutation of the \(x_{i}\) it is necessary and sufficient that the number of irreducible factors of f should be congruent to the degree of f mod 2 (consider a generator \(\sigma\) of G and the decomposition into cycles of the permutation of the \(x_{i}\) induced by \(\sigma\) ).
  \end{enumerate}
\end{enumerate}

\begin{enumerate}
\def\labelenumi{\alph{enumi})}
\setcounter{enumi}{1}
\tightlist
\item
  Suppose further that k is of characteristic \(\neq 2\) . Show that the two above conditions are equivalent to the fact that the discriminant of f should be a square in k (« Stickelberger's theorem »).
\end{enumerate}

\begin{enumerate}
\def\labelenumi{\arabic{enumi})}
\setcounter{enumi}{22}
\tightlist
\item
  Let \(n\) and \(p\) be odd and distinct prime numbers.
\end{enumerate}

\begin{enumerate}
\def\labelenumi{\alph{enumi})}
\item
  Let \(Q \in F_{p}[X]\) be the image of the cyclotomic polynomial \(\Phi_{n}\) . Show that the number i of irreducible factors of Q is even if and only if p is a square mod n.~(Let K be the splitting field of Q. Show that \(\operatorname{Gal}(K/F_{p})\) may be identified with the subgroup of \((\mathbf{F}_{n})^{*}\) generated by the image of p.~Study the parity of the index \(((\mathbf{F}_{n})^{*}: \operatorname{Gal}(K/F_{p}))\) .)
\item
  Show that the discriminant of \(\mathbf{X}^n - 1 \in \mathbf{F}_p[\mathbf{X}]\) is
\end{enumerate}

\[
(- 1) ^ {\frac {n (n - 1)}{2} + n - 1} n ^ {n}.
\]

\begin{enumerate}
\def\labelenumi{\alph{enumi})}
\setcounter{enumi}{2}
\tightlist
\item
  Show that \(i\) is even if and only if \((-1)^{\frac{n(n - 1)}{2}} n\) is a square mod \(p\) (apply Ex. 22 to the polynomial \(\mathbf{X}^n - 1\) ).
\end{enumerate}

\[
\left(\frac {n}{p}\right) = 1
\]

\[
\left(\frac {n}{p}\right) = - 1
\]

\(\left(\frac{n}{p}\right)\left(\frac{p}{n}\right) = (-1)^{\frac{(n - 1)(p - 1)}{4}}\) (« Law of quadratic reciprocity »).
% semantic-fragments: legacy-input-seam
\relax{}
